\documentclass[11pt]{report}

\usepackage[T1]{fontenc}
\usepackage[utf8]{inputenc}
\usepackage{lmodern}
\usepackage[margin=1in]{geometry}
\usepackage{amsmath, amssymb, amsthm, mathtools}
\usepackage{booktabs}
\usepackage{enumitem}
\usepackage{hyperref}
\usepackage{microtype}

\hypersetup{
  colorlinks=true,
  linkcolor=blue,
  urlcolor=blue,
  citecolor=blue,
  plainpages=false,
  pdfpagelabels=true,
  pdftitle={Technical Report: Instance-Specific Portfolio Selection},
  pdfauthor={Christian Speck}
}

\title{Technical Report\\[0.4em]\large Instance-Specific Portfolio Selection for Algorithm Configuration}
\author{Christian Speck}
\date{\today}

\newtheorem{definition}{Definition}[chapter]
\newtheorem{remark}{Remark}[chapter]

\begin{document}

\maketitle
\pagenumbering{roman}
\tableofcontents
\clearpage
\pagenumbering{arabic}

\chapter{Introduction}

The central goal of this project is to study \emph{instance-specific parameter selection} under a restricted action space. We assume that a practitioner has already chosen a finite portfolio of parameter settings for a single underlying algorithm. At decision time, the system observes features of a problem instance and selects one portfolio member. The system does \emph{not} perform online continuous optimization of the parameters. Instead, it performs online \emph{routing} among a fixed set of candidate configurations.

This framing is closely aligned with the motivation of Instance-Specific Algorithm Configuration (ISAC): different regions of the instance space often favor different algorithmic behaviors, and observable features can be used to route a new instance to an appropriate configuration. The main difference in this repository is pedagogical and architectural. Rather than starting with a heavy benchmark and a full clustering pipeline, we begin with a minimal simulation whose structure matches the target decision problem.

\chapter{Problem Formulation}

\section{Instances, Features, and Portfolio Elements}

Let $\mathcal{X}$ denote the space of problem instances. For each instance $x \in \mathcal{X}$, we observe a feature vector
\[
\phi(x) \in \mathbb{R}^d,
\]
where $\phi$ is a feature extractor that is available at decision time.

Let the fixed portfolio of parameter settings be
\[
\Theta = \{\theta^{(1)}, \theta^{(2)}, \dots, \theta^{(K)}\},
\]
where each $\theta^{(k)} \in \mathbb{R}^p$ represents one preselected algorithm configuration. The key modeling choice is that $\Theta$ is finite and fixed before online deployment.

\begin{definition}
A portfolio selector is a mapping
\[
\pi : \mathbb{R}^d \to \{1,2,\dots,K\},
\]
which takes the observable feature vector $\phi(x)$ and returns the index of the portfolio element to execute.
\end{definition}

\section{Performance and Regret}

Let $c(x, \theta^{(k)})$ denote the cost of running configuration $\theta^{(k)}$ on instance $x$. Depending on the domain, this cost may be runtime, loss, error, or a weighted objective combining several metrics. We assume that lower values are better.

The oracle cost for instance $x$ is
\[
c^*(x) = \min_{1 \le k \le K} c(x, \theta^{(k)}).
\]

The \emph{instance regret} of choosing portfolio element $\theta^{(k)}$ on instance $x$ is
\[
r(x, \theta^{(k)}) = c(x, \theta^{(k)}) - c^*(x).
\]

Thus the expected regret of a selector $\pi$ under an instance distribution $\mathcal{D}$ is
\[
\mathbb{E}_{x \sim \mathcal{D}}
\left[
c(x, \theta^{(\pi(\phi(x)))}) - c^*(x)
\right].
\]

The learning objective is to construct $\pi$ so as to minimize this expected regret.

\section{Why This Is Not Reinforcement Learning}

The online decision in this project is a single-shot contextual routing problem. At each timestep, the system observes features and chooses a portfolio member. The environment does not require long-horizon credit assignment, policy improvement through repeated environmental interaction, or continuous online reconfiguration of the parameter vector.

Mathematically, this is much closer to contextual decision making or supervised cost-sensitive classification than to sequential reinforcement learning. The state is summarized by observable features, the action set is finite, and the per-instance loss is immediately measurable once the selected configuration has been evaluated.

\chapter{Synthetic Benchmark Design}

\section{Latent Regimes}

To make the problem nontrivial, the simulation assumes that instances arise from latent regimes
\[
z \in \{1,2,\dots,M\}.
\]
Each regime corresponds to a subset of the instance space for which different parameter settings are advantageous.

For regime $z$, we define:
\begin{itemize}[leftmargin=2em]
\item a feature-space center $\mu_z \in \mathbb{R}^d$,
\item an ideal parameter vector $\theta_z^* \in \mathbb{R}^p$.
\end{itemize}

An instance from regime $z$ is generated by sampling features according to
\[
\phi(x) = \mu_z + \varepsilon_x,
\quad
\varepsilon_x \sim \mathcal{N}(0, \sigma_x^2 I),
\]
followed by an ideal parameter vector
\[
\tilde{\theta}(x) = \operatorname{clip}\bigl(\theta_z^* + g(\phi(x)) + \varepsilon_\theta, 0, 1\bigr),
\quad
\varepsilon_\theta \sim \mathcal{N}(0, \sigma_\theta^2 I),
\]
where $g(\phi(x))$ introduces mild within-regime variation and $\operatorname{clip}$ enforces bounded parameter values.

\section{Cost Model}

Each portfolio element is a fixed vector $\theta^{(k)}$. The benchmark defines a base instance difficulty $b(x)$ and then measures mismatch between the chosen configuration and the latent ideal configuration:
\[
c(x, \theta^{(k)}) = b(x) + \lambda \lVert \theta^{(k)} - \tilde{\theta}(x) \rVert_2^2 + \eta,
\]
where $\lambda > 0$ is a mismatch penalty and $\eta$ is small observational noise.

This cost model yields three important properties:
\begin{enumerate}[leftmargin=2em]
\item Different regimes can favor different portfolio elements.
\item Features are predictive, but not perfectly predictive, of the best configuration.
\item A finite portfolio can be objectively evaluated by regret against the instance-wise best portfolio member.
\end{enumerate}

\section{Current Repository Instantiation}

The current implementation in the repository defines a compact synthetic portfolio benchmark with:
\begin{itemize}[leftmargin=2em]
\item three latent regimes,
\item a small number of observable features,
\item a two-dimensional parameter vector,
\item a four-element fixed portfolio containing conservative, aggressive, balanced, and robust parameter settings.
\end{itemize}

This benchmark lives in \texttt{src/isac/core/portfolio.py}. Its role is to provide a minimal experimental object with the same decision geometry as the intended real problem.

\chapter{Relation to ISAC}

\section{Conceptual Alignment}

ISAC begins from the hypothesis that feature-based structure in the instance space can be exploited to improve algorithm configuration. In broad terms, the logic is:
\begin{enumerate}[leftmargin=2em]
\item represent instances by feature vectors,
\item normalize or preprocess those features,
\item partition the instance space,
\item choose or learn a good configuration for each region,
\item route future instances according to their observed features.
\end{enumerate}

The benchmark in this repository captures the same core idea while simplifying the experimental substrate. Instead of immediately implementing the full offline clustering and per-cluster configuration optimization pipeline, we start from a benchmark where the true performance landscape is controllable and interpretable.

\section{A Minimal Selector Family}

The most direct selector family to study after the benchmark is implemented is a nearest-prototype rule. Suppose we compute centroids $\hat{\mu}_1, \dots, \hat{\mu}_M$ in feature space and estimate the empirically best portfolio member in each region, denoted $k_1^*, \dots, k_M^*$. Then a simple selector has the form
\[
\pi(\phi(x)) = k_{j(x)}^*,
\quad
j(x) = \arg\min_{1 \le m \le M} \lVert \phi(x) - \hat{\mu}_m \rVert_2^2.
\]

This style of rule is attractive because it is transparent, cheap at inference time, and closely related to the classical ISAC intuition.

\chapter{Evaluation Quantities}

\section{Average Cost and Regret}

For a dataset $\{x_i\}_{i=1}^n$, the empirical average cost of a selector $\pi$ is
\[
\hat{C}(\pi) = \frac{1}{n} \sum_{i=1}^n c\bigl(x_i, \theta^{(\pi(\phi(x_i)))}\bigr).
\]
The empirical average regret is
\[
\hat{R}(\pi) = \frac{1}{n} \sum_{i=1}^n
\left(
c\bigl(x_i, \theta^{(\pi(\phi(x_i)))}\bigr) -
\min_{1 \le k \le K} c(x_i, \theta^{(k)})
\right).
\]

\section{Global Fallback Baseline}

A natural baseline is the best \emph{instance-oblivious} portfolio member:
\[
k_{\mathrm{global}}^* =
\arg\min_{1 \le k \le K}
\frac{1}{n} \sum_{i=1}^n c(x_i, \theta^{(k)}).
\]

This baseline is important because it answers a concrete question: how much value do the features provide beyond a single globally fixed parameter setting?

\section{Separability Diagnostics}

Before sophisticated methods are introduced, the benchmark should support basic diagnostics:
\begin{itemize}[leftmargin=2em]
\item the proportion of instances on which each portfolio element is optimal,
\item cluster or regime confusion under simple selectors,
\item regret reduction relative to the global baseline,
\item sensitivity to feature noise and overlap.
\end{itemize}

\chapter{Path to Real Data}

\section{ASlib as the Most Natural Follow-Up}

The most natural real-data path is the Algorithm Selection Library (ASlib). ASlib standardizes scenarios that provide instance features and per-instance performance tables over a finite set of candidate solvers. This is extremely close to the desired project structure.

For the purposes of this repository, one may reinterpret each candidate solver in an ASlib scenario as a \emph{portfolio element}. In some settings that portfolio element is literally a different algorithm. In others, it can be viewed as a stand-in for a precomputed parameterization of one base method. The important structural point is the same: a finite action set is available offline, and the online task is to select among those actions from features.

\section{OpenML-Derived Scenarios}

OpenML-derived scenarios are particularly appealing because they more closely resemble machine-learning model and hyperparameter selection at the dataset level. These scenarios are likely to be a good bridge between the synthetic benchmark and the eventual target application in which one has dataset characteristics and must choose among several fixed configurations.

\chapter{Implementation Notes}

\section{Repository Role of the Report}

This report is intended to function as a living technical chapter. Whenever the repository gains a new benchmark, selector, dataset loader, or evaluation protocol, the report should be updated so that:
\begin{itemize}[leftmargin=2em]
\item the notation remains consistent,
\item the mathematical objective is explicit,
\item the experimental abstractions are documented in prose rather than only in code.
\end{itemize}

\section{Build Process}

The PDF version of this report is generated from this \LaTeX{} source using \texttt{latexmk}. The repository includes a small build script so that the command
\begin{verbatim}
./scripts/build_report.sh
\end{verbatim}
produces an updated PDF in \texttt{docs/report/build/technical\_report.pdf}.

\chapter{Conclusion}

The project begins with a deliberately narrow question: can observable instance characteristics guide the online selection of one configuration from a finite, precomputed parameter portfolio? That question is already rich enough to support meaningful mathematical definitions, interpretable simulations, and rigorous evaluation. By establishing that foundation first, the repository can later expand toward richer ISAC-inspired models without losing clarity about the underlying decision problem.

\end{document}
